\section*{Capítulo \ref{ch:g6:measurement-in-science} --- Medir na ciência: unidades e instrumentos}

\begin{solution}{exo:g6:measurement-in-science:1}
$2$ \unit{m}; \ $8$ \unit{g}; \ $55$ \unit{min}; \ $39$
\unit{\celsius}.
\end{solution}

\begin{solution}{exo:g6:measurement-in-science:2}
\qty{320}{cm}; \ \qty{0.75}{kg}; \ \qty{240}{s}; \ \qty{5.6}{km}.
\end{solution}

\begin{solution}{exo:g6:measurement-in-science:3}
Os degraus dele são de $60$, e não de $10$, então a escada decimal
mente: $\qty{3.5}{min} = 3 \times 60 + 30 = \qty{210}{s}$, enquanto a
escada teria dado erradamente \qty{350}{s}.
\end{solution}

\begin{solution}{exo:g6:measurement-in-science:4}
Cada passinho: $200 \div 5 = \qty{40}{mL}$. Dois passinhos acima de
$300$: $300 + 80 = \qty{380}{mL}$.
\end{solution}

\begin{solution}{exo:g6:measurement-in-science:5}
Fermento: a balança de cozinha (passinhos de um \omterm{def:g3:measuring-mass:units}{grama} --- a balança
de banheiro nem notaria uma pitada). Batimentos: um cronômetro
(passinhos de segundo e menores). Largura da página: uma régua
(passinhos de milímetro). Pátio da escola: uma trena comprida ou uma
roda de medição (alcance na escala de \omterm{def:g2:measuring-length:units}{metros}).
\end{solution}

\begin{solution}{exo:g6:measurement-in-science:6}
O passo do julgamento: um lápis é estimado em menos de \qty{20}{cm},
então \qty{17.4}{m} reprova na hora no teste de bom senso. Ela leu
\omterm{def:g2:measuring-length:units}{centímetros} e escreveu \omterm{def:g2:measuring-length:units}{metros}: o valor verdadeiro é \qty{17.4}{cm}.
\end{solution}

\begin{solution}{exo:g6:measurement-in-science:7}
Honesto: ``entre $18$ e \qty{19}{\celsius}, mais perto de $19$'' (ou
``cerca de \qty{18.8}{\celsius}''). Desonesto:
``\qty{18.7643}{\celsius}'' --- uma precisão que as marcas não podem
prometer.
\end{solution}

\begin{solution}{exo:g6:measurement-in-science:8}
Ninguém errou: os dedos começam e param um fiozinho antes ou depois,
e uma dispersão pequena é a respiração normal de uma medida. Relato
sábio: as leituras se agrupam em torno de \qty{6.2}{s} --- ``cerca de
\qty{6.2}{s}''.
\end{solution}

\begin{solution}{exo:g6:measurement-in-science:9}
Um quarto de \qty{450}{g}: cerca de \qty{112}{g} (\qty{113}{g}
também serve). A tarefa mostra por que \omterm{def:g6:measurement-in-science:measuring}{unidades} compartilhadas
importam: receitas escritas nas \omterm{def:g6:measurement-in-science:measuring}{unidades} de um país precisam ser
traduzidas antes que os instrumentos de outro país possam obedecer.
\end{solution}

\begin{solution}{exo:g6:measurement-in-science:10}
O selo vai da marca $1$ até a marca $3.4$: $3.4 - 1 = \qty{2.4}{cm}$.
O hábito: nunca confiar na borda de um instrumento --- comece de uma
marca escolhida e subtraia, e as réguas gastas perdem o poder de
enganar.
\end{solution}

\begin{solution}{exo:g6:measurement-in-science:11}
As duas leituras se espalham por \qty{0.10}{m}, então o último
algarismo de qualquer uma delas não é confiável; elas se agrupam em
torno de $18.3$. Escrever ``\qty{18.3}{m}'' guarda todos os
algarismos que a trena pode prometer com honestidade, e nada além
disso.
\end{solution}

\begin{solution}{exo:g6:measurement-in-science:12}
Desastres: os ``pés'' de duas cidades nunca coincidem, então vigas
cortadas numa cidade nunca encaixam nas casas de outra; e, quando o
filho herda o trono, a própria \omterm{def:g6:measurement-in-science:measuring}{unidade} muda de comprimento --- toda
medida do reino vence junto com o rei. A barra de metal resolveu os
dois: um comprimento imutável, copiado para todas as cidades,
sobrevivendo a todos os reis. O defeito dela: um objeto único pode
ser arranhado, roubado ou alterado devagarinho, e o mundo tinha de
peregrinar até o cofre dele para conferir as cópias --- até que as
\omterm{def:g6:measurement-in-science:measuring}{unidades} foram reconstruídas sobre as constantes da natureza, que não
se desgastam e estão de plantão em toda parte ao mesmo tempo.
\end{solution}

\begin{solution}{pb:g6:measurement-in-science:1}
\textbf{1.} Toda pesagem soma a carga ao que o ponteiro já mostrava;
um começo diferente de zero envenena todas as leituras do dia --- a
conferência do começo justo tem de vir primeiro.
\textbf{2.} $500 + 200 = \qty{700}{g}$.
\textbf{3.} Ela marca \qty{20}{g} a menos. Despejando farinha até o
ponteiro dizer \qty{700}{g}, o padeiro despejou na verdade
\qty{720}{g}: a balança faz o padeiro entregar mais mercadoria do que
o freguês paga --- ela favorece os fregueses, e o padeiro perde em
cada venda.
\textbf{4.} $500 + 200 + 50 = \qty{750}{g}$.
\textbf{5.} Não: um selo de \qty{2}{g} some abaixo dos passinhos de
\qty{5}{g} da balança --- está fora do território dela.
\textbf{6.} $200 \div 4 = \qty{50}{mL}$ por passinho.
\textbf{7.} A água está em $600 + 50 = \qty{650}{mL}$; ``cerca de
\qty{700}{mL}'' arredonda um passinho inteiro para longe --- vago
demais para um instrumento tão bom.
\textbf{8.} O olho na mesma altura da marca: ler de baixo desloca a
linha aparente da água --- a regra de olhar bem de frente foi
quebrada.
\textbf{9.} Errada por \qty{10}{mL} --- melhor (menor) que o próprio
passinho de \qty{50}{mL}: aceitável para uma jarra de escola.
\textbf{10.} Adiantado, em $2$ minutos.
\textbf{11.} $2 \times 7 = 14$ minutos adiantado depois de uma
semana.
\textbf{12.} Primeiro relógio: atrasar $2$ min; segundo: adiantar $2$
min; terceiro: nada --- ele está certo. (Tirar a média dos três não
``consertaria'' nada e falsificaria o único relógio honesto.)
\textbf{13.} Por exemplo: \emph{Zero}: confira se o instrumento vazio
marca zero antes de confiar em qualquer coisa que ele diga.
\emph{Leitura}: descubra quanto vale o passinho e leia bem de frente,
com modéstia honesta entre as marcas. \emph{\omterm{def:g6:measurement-in-science:measuring}{Unidades}}: um número sem
a \omterm{def:g6:measurement-in-science:measuring}{unidade} dele não é uma medida --- escreva a \omterm{def:g6:measurement-in-science:measuring}{unidade}, sempre.
\end{solution}
